\documentclass[12pt, a4paper]{article}
\usepackage{amsmath, amssymb, amsthm}
\usepackage{geometry}
\geometry{a4paper, margin=1in}
\usepackage{hyperref}
\usepackage{graphicx}
\usepackage{fancyhdr}
\usepackage{cleveref}
\usepackage{tcolorbox}
\usepackage{enumitem}
\usepackage{datetime}

% Cross-referencing improvements
\usepackage{xr}
\usepackage{zref-xr}
\usepackage{zref-user}

% Document Information
\title{The Meta-Holo-Morphic Framework: A Multidisciplinary Approach to the\\Millennium Prize Problems and the Grand Unified Theory of Physics}
\author{Thomas Birnie}
\date{\today}

% Header and Footer Setup
\pagestyle{fancy}
\fancyhead[L]{Millennium Prize Problems}
\fancyhead[R]{\thepage}
\setlength{\headheight}{14.5pt} % Fix for fancyhdr warning

% Graphics Path
\graphicspath{{./images/}}

% Theorem, Lemma, Definition, Corollary, and Proposition Environments
\newtheorem{theorem}{Theorem}[section]
\newtheorem{lemma}[theorem]{Lemma}
\newtheoremstyle{definitionstyle}
  {0.5em}   % Space above
  {0.5em}   % Space below
  {}        % Body font
  {}        % Indent amount
  {\bfseries} % Theorem head font
  {.}        % Punctuation after theorem head
  {.5em}     % Space after theorem head
  {}         % Theorem head spec
\theoremstyle{definitionstyle}
\newtheorem{definition}[theorem]{Definition}
\newtheorem{corollary}[theorem]{Corollary}
\newtheorem{proposition}[theorem]{Proposition}

% Enhanced mathematical environments
\newtheorem{conjecture}[theorem]{Conjecture}
\newtheorem{algorithm}{Algorithm}[section]
\newtheorem{observation}[theorem]{Observation}

% Conceptbox for important highlights
\newenvironment{conceptbox}[1]
{\begin{tcolorbox}[title=#1]}
{\end{tcolorbox}}

% Definitions and Notation Section Command
\newcommand{\notationsection}{
    \section*{Definitions and Notation}
    \addcontentsline{toc}{section}{Definitions and Notation}
    % Insert definitions and notations here
}

% Proof Formatting
\renewenvironment{proof}[1][Proof]{\noindent\textbf{#1.} }{\hfill $\square$\vspace{0.5em}}

% Mathematical notation macros
\newcommand{\N}{\mathbb{N}}
\newcommand{\Z}{\mathbb{Z}}
\newcommand{\Q}{\mathbb{Q}}
\newcommand{\R}{\mathbb{R}}
\newcommand{\C}{\mathbb{C}}

% Proof strategy environment
\newenvironment{proofstrategy}
  {\begin{tcolorbox}[title=Proof Strategy]%
   \begin{enumerate}[label=\arabic*.]}
  {\end{enumerate}\end{tcolorbox}}

% Mathematical diagram support
\usepackage{tikz}
\usetikzlibrary{cd,matrix,arrows,decorations.pathmorphing}

\begin{document}

\maketitle

\tableofcontents
\newpage

% Include generated sections
\input{generated_sections.tex}

% Introduction
\section{Introduction}
\subsection{Purpose and Scope}
This document provides a comprehensive exploration of the Millennium Prize Problems (MPPs), integrating foundational mathematical tools, modern physics theories, and philosophical perspectives. The content is structured to offer a unified approach to understanding these profound questions, intended for submission to the Clay Mathematics Institute. Each section aims to build upon the previous one, creating a cohesive and in-depth journey through some of the most significant challenges in modern mathematics and physics.

\subsection{Motivation}
The purpose of this document is to outline the criteria set by the Clay Mathematics Institute for proofs, setting the foundation for the approach to each problem. It establishes the overarching goals of this document and how it seeks to contribute to the current understanding of these deep questions.

\subsection{Visual Aids}
\begin{figure}[htbp]
   \centering
   \includegraphics[width=\textwidth]{intro_visual.png}
   \caption{Overview of CMI standards and Millennium Prize Problems}
   \label{fig:intro_visual}
\end{figure}

% Definitions and Notation
\notationsection

% 1. Birch and Swinnerton-Dyer Conjecture
\section{Birch and Swinnerton-Dyer Conjecture}
\subsection{Elliptic Curves and L-Functions}
\begin{definition}[Elliptic Curve]
An elliptic curve $E$ over $\mathbb{Q}$ is a smooth, projective algebraic curve of genus one, with a specified point $O$, defined over the field of rational numbers.
\end{definition}

\begin{definition}[L-function]
The L-function $L(E,s)$ of an elliptic curve $E$ over $\mathbb{Q}$ is defined by an Euler product over all primes $p$:
\begin{equation}
L(E,s) = \prod_{p\,\text{prime}} L_p(p^{-s})^{-1},
\end{equation}
where $L_p(p^{-s})$ encodes information about the number of points on $E$ modulo $p$.
\end{definition}

\subsection{Rational Points and Rank of Elliptic Curves}
\begin{definition}[Rank of an Elliptic Curve]
The rank of $E(\mathbb{Q})$ is the number of independent infinite order points in the group of rational points on $E$.
\end{definition}

\subsection{Theorem Statement}
\begin{theorem}[Birch and Swinnerton-Dyer Conjecture]
\label{thm:bsd}
The rank of an elliptic curve $E$ over $\mathbb{Q}$ is equal to the order of the zero of $L(E,s)$ at $s=1$.
\end{theorem}

\subsection{Tools and Techniques}
% Discuss computational methods and modern algorithms

\subsection{Case Studies}
% Provide analysis of specific elliptic curves

\subsection{Connections to Other Problems}
% Highlight connections to other MPPs

% 2. Hodge Conjecture
\section{Hodge Conjecture}
\subsection{Algebraic Geometry Framework}
\begin{definition}[Kähler Manifold]
A Kähler manifold is a complex manifold equipped with a Hermitian metric whose associated two-form is closed.
\end{definition}

\begin{definition}[Hodge Decomposition]
For a compact Kähler manifold $X$, the cohomology group $H^k(X, \mathbb{C})$ admits a decomposition:
\[
H^k(X, \mathbb{C}) = \bigoplus_{p+q=k} H^{p,q}(X),
\]
where $H^{p,q}(X)$ are the spaces of harmonic forms of type $(p,q)$.
\end{definition}

\subsection{Hodge Structures and Decomposition}
% Discuss the intricate structures and their significance

\subsection{Tools and Techniques}
% Explain algebraic topology tools used

\subsection{Illustrative Examples}
% Provide case studies involving complex projective varieties

\subsection{Visual Aids}
% Include diagrams of cohomology classes and algebraic cycles

% 3. Riemann Zeta Theorem
\section{Riemann Zeta Theorem}
\subsection{Statement of the Hypothesis}
\begin{theorem}[Riemann Hypothesis]
\label{thm:riemann}
All non-trivial zeros of the Riemann zeta function $\zeta(s)$ have real part $\frac{1}{2}$; that is, they lie on the critical line $\Re(s) = \frac{1}{2}$ in the complex plane.
\end{theorem}

\subsection{Prime Distribution and Implications}
% Discuss how this relates to the distribution of prime numbers

\subsection{Analytic Techniques}
% Introduce complex analysis techniques used in studying $\zeta(s)$

\subsection{Connections to Quantum Chaos}
% Explore analogies between zeta function behavior and quantum systems

\subsection{Visual Aids}
\begin{figure}[htbp]
   \centering
   \includegraphics[width=\textwidth]{riemann_zeros.png}
   \caption{Visualization of the non-trivial zeros of the Riemann zeta function on the critical line}
   \label{fig:riemann_zeros}
\end{figure}

% 4. Navier-Stokes Equations
\section{Navier-Stokes Equations}
\subsection{Introduction to Fluid Dynamics}
\begin{definition}[Navier-Stokes Equations]
The Navier-Stokes equations for incompressible fluid flow are:
\begin{align}
\frac{\partial \mathbf{u}}{\partial t} + (\mathbf{u} \cdot \nabla)\mathbf{u} &= -\nabla p + \nu \Delta \mathbf{u} + \mathbf{f}, \\
\nabla \cdot \mathbf{u} &= 0,
\end{align}
where $\mathbf{u}$ is the velocity field, $p$ is pressure, $\nu$ is viscosity, and $\mathbf{f}$ represents external forces.
\end{definition}

\subsection{Analytical Approach to Regularity}
% Introduce functional analysis tools for understanding solution regularity

\subsection{Bounding Energy and Controlling Blow-Ups}
% Discuss techniques for bounding energy and preventing singularities

\subsection{Application of Morphic Functions}
% Explore the use of morphic functions in complex fluid flows

\subsection{Case Studies and Visual Aids}
% Provide real-world examples and include fluid flow diagrams

% 5. Quantum Holo-Morphic Framework
\section{Quantum Holo-Morphic Framework}
\subsection{Introduction to Holography}
\begin{definition}[Holographic Principle]
A principle stating that the description of a volume of space can be thought of as encoded on a lower-dimensional boundary to the region.
\end{definition}

\subsection{Applications to Black Holes}
% Explore black hole thermodynamics and thought experiments

\subsection{Complex Time and Unification}
% Introduce the concept of complex time in unifying theories

\subsection{Visual Aids}
% Include diagrams of black holes and holographic projections

% 6. Yang-Mills Existence and Mass Gap
\section{Yang-Mills Existence and Mass Gap}
\subsection{Introduction to Yang-Mills Theory}
\begin{definition}[Yang-Mills Equations]
Non-abelian generalization of Maxwell's equations, fundamental in quantum field theory.
\end{definition}

\subsection{Compactification and Gauge Symmetries}
% Introduce gauge symmetries and discuss mass gaps

\subsection{Mathematical Proofs and Tools}
% Provide techniques to address the mass gap problem

\subsection{Visual Aids}
% Diagram of field configurations and compactification

% 7. Theory of Everything
\section{Theory of Everything}
\subsection{Unification of Forces}
% Describe fundamental forces and efforts towards unification

\subsection{Integration with Quantum Gravity}
% Explore models of quantum gravity

\subsection{Symmetry Breaking and Vacuum Structure}
% Analyze vacuum states and symmetry breaking

\subsection{Visual Aids}
% Diagrams illustrating unification of forces

% 8. P vs NP Problem
\section{P vs NP Problem}
\subsection{Statement of the Problem}
\begin{theorem}[P vs NP]
The question of whether every problem whose solution can be quickly verified can also be solved quickly.
\end{theorem}

\subsection{Recursive Reduction Framework}
% Introduce symmetry-breaking, recursive decomposition, and canonical forms

\subsection{Formalizing Lemmas}
% Present and prove lemmas for complexity reduction

\subsection{Agent Models and Practical Implications}
% Introduce two-agent model for solving NP problems

\subsection{Visual Aids}
\begin{figure}[htbp]
   \centering
   \includegraphics[width=\textwidth]{pnp_visual.png}
   \caption{Visualization of computational complexity classes}
   \label{fig:pnp_visual}
\end{figure}

% 9. Conclusion
\section{Conclusion}
\subsection{Summary of Key Contributions}
% Provide an overview of contributions to solving each MPP

\subsection{Implications for Future Research}
% Discuss broader implications for mathematics and physics

\subsection{Final Thoughts}
% Reflect on the unified approach to mathematical reality

% Appendices
\appendix
\section{Technical Background}
% Supplementary material on tools used

\section{Glossary}
% Glossary of terms

\section{Additional Visual Aids}
% Additional diagrams, graphs, and illustrations

% Bibliography
\bibliographystyle{plain}
\bibliography{references}

\end{document}
